\documentclass[12pt, a4paper]{article}

\usepackage[utf8]{inputenc}
\usepackage[T1]{fontenc}
\usepackage[margin=2.5cm, headheight=15pt]{geometry}
\usepackage{titlesec}
\usepackage{parskip}
\usepackage{setspace}
\setstretch{1.15}
\usepackage{hyperref}
\usepackage{listings}
\usepackage{xcolor}
\usepackage{booktabs}
\usepackage{enumitem}
\usepackage{fancyhdr}
\usepackage{graphicx}

% Code listing style
\lstdefinestyle{cstyle}{
    language=C,
    basicstyle=\ttfamily\small,
    keywordstyle=\color{blue}\bfseries,
    commentstyle=\color{gray}\itshape,
    stringstyle=\color{red},
    numberstyle=\tiny\color{gray},
    numbers=left,
    stepnumber=1,
    frame=single,
    breaklines=true,
    backgroundcolor=\color{gray!8},
    captionpos=b,
    tabsize=4
}

\lstdefinestyle{bashstyle}{
    language=bash,
    basicstyle=\ttfamily\small,
    keywordstyle=\color{blue},
    commentstyle=\color{gray}\itshape,
    frame=single,
    breaklines=true,
    backgroundcolor=\color{gray!8},
}

% Header/Footer
\pagestyle{fancy}
\fancyhf{}
\rhead{CSE344 -- Homework \#1}
\lhead{Halit Batuhan ASLAN -- 220104004003}
\rfoot{\thepage}

% Section formatting
\titleformat{\section}{\large\bfseries}{}{0em}{}[\titlerule]
\titlespacing{\section}{0pt}{18pt}{10pt}
\titleformat{\subsection}{\normalsize\bfseries}{}{0em}{}
\titlespacing{\subsection}{0pt}{14pt}{6pt}

\hypersetup{
    colorlinks=true,
    linkcolor=blue,
    urlcolor=blue
}

\begin{document}

% ─── Title Page ───────────────────────────────────────────────────────────────
\begin{titlepage}
    \centering
    \vspace*{2cm}

    {\LARGE \textbf{Gebze Technical University}} \\[0.5cm]
    {\large Department of Computer Engineering} \\[2cm]

    \rule{\linewidth}{0.5mm} \\[0.4cm]
    {\huge \textbf{CSE344 -- System Programming}} \\[0.3cm]
    {\Large \textbf{Homework \#1: Advanced File Search Program}} \\[0.1cm]
    \rule{\linewidth}{0.5mm} \\[2cm]

    {\large \textbf{Student:} Halit Batuhan ASLAN} \\[0.3cm]
    {\large \textbf{Student Number:} 220104004003} \\[0.3cm]
    {\large \textbf{Date:} March 2026} \\[3cm]

    \vfill
    {\small Submitted for CSE344 -- System Programming, Spring 2026}
\end{titlepage}

% ─── Table of Contents ────────────────────────────────────────────────────────
\tableofcontents
\newpage

% ─── 1. Introduction ──────────────────────────────────────────────────────────
\section{Introduction}

The objective of this homework is to implement an advanced file search program for
POSIX-compatible operating systems. The program, named \texttt{myFind}, recursively
traverses a given directory tree and searches for files and directories that satisfy
a set of criteria specified by the user via command-line arguments.

The program supports the following search criteria:

\begin{itemize}[noitemsep]
    \item \textbf{-f}: Filename pattern matching (case-insensitive, with \texttt{+} regex operator support)
    \item \textbf{-b}: File size in bytes (exact match)
    \item \textbf{-t}: File type (regular file, directory, symbolic link, socket, block device, character device, pipe)
    \item \textbf{-p}: File permissions as a 9-character string (e.g. \texttt{rwxr-xr-x})
    \item \textbf{-l}: Number of hard links (exact match)
\end{itemize}

The mandatory argument \textbf{-w} specifies the root directory for the recursive search.
At least one search criterion must be provided. Results are printed in a tree-formatted
output to stdout, while all error messages go to stderr. The program also handles
CTRL-C (SIGINT) gracefully by printing an informational message and exiting cleanly.

% ─── 2. Design Decisions ──────────────────────────────────────────────────────
\section{Design Decisions}

\subsection{Modular Structure}

The program is divided into seven modules, each with a single, well-defined responsibility.
This separation of concerns makes the code easier to understand, test, and maintain.
Each module consists of a header file (\texttt{.h}) and an implementation file (\texttt{.c}):

\begin{center}
\begin{tabular}{lll}
\toprule
\textbf{Module} & \textbf{Files} & \textbf{Responsibility} \\
\midrule
Argument Parser  & \texttt{parse\_argv.h/.c}      & Parse and validate command-line arguments \\
File Checker     & \texttt{check\_files.h/.c}     & Check if a file matches all criteria \\
Directory Search & \texttt{searching\_files.h/.c} & Recursive directory traversal \\
Tree Printer     & \texttt{print\_tree.h/.c}      & Format and print the tree output \\
Regex Engine     & \texttt{regex.h/.c}            & Pattern matching with \texttt{+} operator \\
Signal Handler   & \texttt{stopping\_handler.h/.c}& Handle SIGINT (CTRL-C) gracefully \\
Entry Point      & \texttt{main.c}                & Program orchestration \\
\bottomrule
\end{tabular}
\end{center}

\subsection{Use of \texttt{lstat()} Instead of \texttt{stat()}}

All file metadata retrieval uses \texttt{lstat()} rather than \texttt{stat()}. This is a deliberate
design decision: \texttt{stat()} follows symbolic links and returns the metadata of the target
file, whereas \texttt{lstat()} examines the symbolic link itself. This is required for the
\texttt{-t l} criterion to correctly identify symbolic links rather than the files they point to.

\subsection{Sentinel Values for Unused Criteria}

The \texttt{Search\_criteria} struct uses sentinel values to distinguish between ``not provided''
and a valid value of zero:

\begin{itemize}[noitemsep]
    \item Pointer fields (\texttt{file\_name}, \texttt{permissions}): \texttt{NULL}
    \item \texttt{file\_size}: \texttt{-1} (since valid sizes are $\geq 0$)
    \item \texttt{file\_type}: \texttt{'\textbackslash0'} (null character)
    \item \texttt{link\_count}: \texttt{-1} (since valid link counts are $\geq 1$)
\end{itemize}

In \texttt{check\_matches()}, each criterion is checked only if its field is not set to its
sentinel value. This means multiple criteria form an AND condition: all provided criteria
must be satisfied simultaneously.

\subsection{Custom \texttt{+} Regex Operator}

The \texttt{+} operator is implemented from scratch without any external library.
The algorithm uses two nested loops:

\begin{itemize}[noitemsep]
    \item The \textbf{outer loop} tries every position in the string as a potential starting point.
    \item The \textbf{inner loop} matches the pattern character by character, advancing an independent
          offset variable \texttt{k} to handle the \texttt{+} operator correctly.
    \item When \texttt{pattern[j] == '+'}, the preceding character is greedily consumed: the
          algorithm advances \texttt{k} as long as the current string character matches
          \texttt{pattern[j-1]}. The bounds check \texttt{i + k < len\_of\_string} prevents
          buffer overflows.
\end{itemize}

The matching is case-insensitive: both characters are converted to lowercase via
\texttt{tolower()} before comparison.

\subsection{SIGINT Handling with \texttt{volatile sig\_atomic\_t}}

Rather than calling \texttt{exit()} directly from the signal handler, the program uses a
shared flag \texttt{continue\_running} of type \texttt{volatile sig\_atomic\_t}:

\begin{itemize}[noitemsep]
    \item \texttt{volatile} prevents the compiler from caching the value in a register,
          ensuring every read fetches the current value from memory.
    \item \texttt{sig\_atomic\_t} guarantees that reads and writes to this variable are
          atomic, which is a strict requirement for variables modified inside signal handlers.
    \item The flag is checked at the beginning of every iteration in \texttt{search\_directory()}.
          When set to \texttt{0}, the function closes the open directory and returns immediately,
          unwinding all recursive calls cleanly.
\end{itemize}

\texttt{sigaction()} is used instead of the legacy \texttt{signal()} function because
\texttt{sigaction()} has well-defined, consistent behavior across all POSIX systems.

\subsection{Use of \texttt{static} for Internal Helper Functions}

All helper functions that are not part of the public interface of a module are declared
\texttt{static}. This restricts their visibility to the translation unit (the \texttt{.c} file)
in which they are defined. The benefits are twofold: it prevents accidental use from other
modules, and it allows the compiler to optimize more aggressively since it knows the function
cannot be called from outside. Examples include \texttt{char\_equal()} in \texttt{regex.c},
\texttt{check\_type()} and \texttt{check\_permissions()} in \texttt{check\_files.c}, and
\texttt{is\_dot\_or\_dotdot()} and \texttt{build\_path()} in \texttt{searching\_files.c}.

\subsection{AND Logic for Multiple Criteria}

When the user provides more than one search criterion, the program requires all of them
to be satisfied simultaneously (AND logic). This is the most intuitive and useful behavior:
for example, \texttt{-f passwd -t f} means ``a regular file named passwd'', not ``anything
named passwd OR any regular file''. The implementation naturally achieves this by returning
\texttt{0} from \texttt{check\_matches()} as soon as any single criterion fails, without
evaluating the remaining ones (short-circuit evaluation).

\subsection{\texttt{MAX\_PATH\_LENGTH = 4096}}

The buffer used to construct child paths in \texttt{search\_directory()} is sized at
\texttt{4096} bytes. This value corresponds to \texttt{PATH\_MAX} as defined in
\texttt{<linux/limits.h>} on Linux systems, which is the maximum length of a full
absolute path. Using \texttt{snprintf()} with this limit ensures that path construction
never overflows the buffer, even for deeply nested directory trees with long names.

\subsection{Directory Traversal Strategy}

For every entry in a directory, the program:
\begin{enumerate}[noitemsep]
    \item Skips \texttt{.} and \texttt{..} to prevent infinite recursion.
    \item Builds the full child path using \texttt{snprintf()} with a buffer of \texttt{MAX\_PATH\_LENGTH = 4096} bytes.
    \item Calls \texttt{lstat()} to retrieve file metadata.
    \item If the entry is a directory: checks it against criteria, prints if matched, then recurses.
    \item If the entry is any other type: checks it against criteria and prints if matched.
\end{enumerate}

This strategy ensures the entire subtree is always traversed regardless of whether
a directory itself matches the search criteria.

% ─── 3. Modules ───────────────────────────────────────────────────────────────
\section{Modules and Their Explanations}

\subsection{\texttt{parse\_argv} -- Argument Parser}

This module uses \texttt{getopt()} to parse command-line arguments in the format
\texttt{"w:f:b:t:p:l:"}. Each option is validated immediately:

\begin{itemize}[noitemsep]
    \item \textbf{-w}: Stored as a pointer into \texttt{argv} (no copy needed).
    \item \textbf{-f}: Stored as-is; pattern matching is delegated to the regex module.
    \item \textbf{-b}: Converted with \texttt{atol()}; rejected if negative.
    \item \textbf{-t}: Validated against the seven POSIX file type characters.
    \item \textbf{-p}: Validated as a 9-character permission string (\texttt{[rwx-][rwx-][rwx-]}).
    \item \textbf{-l}: Converted with \texttt{atoi()}.
\end{itemize}

If \texttt{-w} is missing, or if no search criterion is provided, the program prints a
usage guide to stderr and exits with code 1.

\subsection{\texttt{check\_files} -- File Criterion Checker}

The \texttt{check\_matches()} function applies all criteria to a single file in sequence.
It returns \texttt{0} (no match) as soon as any criterion fails, avoiding unnecessary
checks. The function uses helper functions \texttt{check\_type()} and
\texttt{check\_permissions()} which are declared \texttt{static} to limit their scope
to this translation unit.

\subsection{\texttt{searching\_files} -- Recursive Directory Traversal}

The \texttt{search\_directory()} function uses \texttt{opendir()}/\texttt{readdir()}/\texttt{closedir()}
for directory traversal. The \texttt{found\_any} pointer is shared across all recursive
calls so that \texttt{main()} can determine whether to print \texttt{"No file found"}.

\subsection{\texttt{print\_tree} -- Tree Formatter}

Two functions handle output formatting:
\begin{itemize}[noitemsep]
    \item \texttt{print\_root()}: Prints the root directory path with no prefix.
    \item \texttt{print\_leaf()}: Prints \texttt{|} followed by \texttt{level * HYPHEN\_PER\_LEVEL} hyphens and the entry name.
\end{itemize}

\subsection{\texttt{regex} -- Pattern Matching Engine}

Implements case-insensitive substring search with the \texttt{+} operator. No external
regex library is used. The \texttt{char\_equal()} helper uses \texttt{tolower()} with an
\texttt{(unsigned char)} cast to avoid undefined behavior on signed character systems.

\subsection{\texttt{stopping\_handler} -- Signal Handler}

Registers \texttt{handle\_sigint()} as the SIGINT handler via \texttt{sigaction()}.
The handler sets \texttt{continue\_running = 0} and prints an exit message to stderr.
The \texttt{(void)signum} suppresses the unused parameter warning from \texttt{-Wall -Wextra}.

% ─── 4. File Structure & Build ────────────────────────────────────────────────
\section{File Structure and Build Instructions}

\subsection{Project File Structure}

The submission directory is organized as follows:

\begin{center}
\begin{tabular}{ll}
\toprule
\textbf{File} & \textbf{Purpose} \\
\midrule
\texttt{main.c}                & Entry point, program orchestration \\
\texttt{parse\_argv.h/.c}      & Command-line argument parsing and validation \\
\texttt{check\_files.h/.c}     & File criteria checking \\
\texttt{searching\_files.h/.c} & Recursive directory traversal \\
\texttt{print\_tree.h/.c}      & Tree-formatted output \\
\texttt{regex.h/.c}            & Custom \texttt{+} regex pattern matching \\
\texttt{stopping\_handler.h/.c}& SIGINT signal handling \\
\texttt{Makefile}              & Build configuration \\
\texttt{report.pdf}            & This report \\
\texttt{report.tex}            & LaTeX source of this report \\
\bottomrule
\end{tabular}
\end{center}

\subsection{Build Instructions}

The program is compiled using the provided \texttt{Makefile}. The following commands
are available:

\begin{itemize}[noitemsep]
    \item \texttt{make} — compiles all source files and produces the executable \texttt{HBA\_file\_search}
    \item \texttt{make clean} — removes all generated object files and the executable
\end{itemize}

The compiler flags used are \texttt{-Wall -Wextra -g}, which enable all standard and
extra warnings and include debug symbols. The program compiles with zero warnings.

\subsection{Usage Examples}

\begin{lstlisting}[style=bashstyle]
# Search for files named "passwd" under /etc
./HBA_file_search -w /etc -f "passwd"

# Search for regular files of exactly 0 bytes
./HBA_file_search -w /tmp -b 0 -t f

# Search for directories with rwxr-xr-x permissions
./HBA_file_search -w /usr -t d -p "rwxr-xr-x"

# Search for symbolic links named "python"
./HBA_file_search -w /usr -f "python" -t l

# Search using the + regex operator
./HBA_file_search -w /etc -f "pas+wd"
\end{lstlisting}

% ─── 5. Requirements Met ──────────────────────────────────────────────────────
\section{Requirements Achieved}

All requirements specified in the homework description have been implemented and tested:

\begin{itemize}[noitemsep]
    \item \textbf{Recursive traversal}: The entire subtree under \texttt{-w} is always traversed.
    \item \textbf{-f criterion}: Case-insensitive filename search with \texttt{+} operator, implemented from scratch.
    \item \textbf{-b criterion}: Exact file size match using \texttt{st\_size} from \texttt{lstat()}.
    \item \textbf{-t criterion}: All seven POSIX file types supported using POSIX macros (\texttt{S\_ISREG}, \texttt{S\_ISDIR}, etc.).
    \item \textbf{-p criterion}: 9-character permission string comparison using bitwise checks on \texttt{st\_mode}.
    \item \textbf{-l criterion}: Hard link count match using \texttt{st\_nlink}.
    \item \textbf{Multiple criteria}: Any combination of criteria is supported (AND logic).
    \item \textbf{Tree output}: Formatted with \texttt{|} and hyphens per depth level.
    \item \textbf{"No file found"}: Printed when no match is found across the entire tree.
    \item \textbf{Error messages to stderr}: All system call failures and invalid arguments are reported to stderr.
    \item \textbf{CTRL-C handling}: Graceful exit with informational message via \texttt{sigaction()} and \texttt{volatile sig\_atomic\_t}.
    \item \textbf{No memory leaks}: Verified with Valgrind (\texttt{All heap blocks were freed -- no leaks are possible}).
    \item \textbf{Usage information}: Printed when mandatory or at-least-one arguments are missing.
    \item \textbf{Makefile}: Provided, compiles with \texttt{-Wall -Wextra -g}, does not run the program.
\end{itemize}

% ─── 6. Testing ───────────────────────────────────────────────────────────────
\section{Testing}

\subsection{Automated Test Suite}

The program was tested using a shell script (\texttt{test.sh}) containing over 45
test cases organized into groups, each targeting a specific feature:

\begin{center}
\begin{tabular}{lll}
\toprule
\textbf{Group} & \textbf{Feature Tested} & \textbf{Result} \\
\midrule
Group 1 & Mandatory \texttt{-w} argument, missing/invalid args & Pass \\
Group 2 & Filename search with \texttt{-f} (exact, case, partial, \texttt{+}) & Pass \\
Group 3 & File type search with \texttt{-t} (f, d, l, s, b, c, p) & Pass \\
Group 4 & File size search with \texttt{-b} (zero, nonzero, nonexistent) & Pass \\
Group 5 & Permission search with \texttt{-p} & Pass \\
Group 6 & Hard link count search with \texttt{-l} & Pass \\
Group 7 & Combined criteria (multiple flags together) & Pass \\
Group 8 & Edge cases (empty directory, no matches) & Pass \\
\bottomrule
\end{tabular}
\end{center}

\subsection{Memory Leak Check with Valgrind}

The program was verified to have zero memory leaks using Valgrind:

\begin{lstlisting}[style=bashstyle]
valgrind --leak-check=full ./HBA_file_search -w /etc -f "passwd"
\end{lstlisting}

Valgrind output confirmed:

\begin{lstlisting}[style=bashstyle]
All heap blocks were freed -- no leaks are possible
ERROR SUMMARY: 0 errors from 0 contexts
\end{lstlisting}

Since the program does not use \texttt{malloc()} or \texttt{free()} directly — all
string pointers point into \texttt{argv} which is managed by the OS — there is no
heap memory to leak. The Valgrind result confirms no indirect leaks either.

\subsection{SIGINT (CTRL-C) Handling Test}

The signal handler was tested by running the program on a large directory tree and
pressing CTRL-C mid-traversal:

\begin{lstlisting}[style=bashstyle]
./HBA_file_search -w / -t f
# Press CTRL-C during execution
\end{lstlisting}

Expected and observed output on stderr:

\begin{lstlisting}[style=bashstyle]
Received SIGINT, exitting program ...
\end{lstlisting}

The program stops traversal cleanly, closes all open directory handles, and exits
without leaving any zombie processes or resource leaks.

\subsection{Edge Cases Tested}

\begin{itemize}[noitemsep]
    \item \textbf{Empty directory}: When the target directory contains no files, the program
          prints only the root path and \texttt{No file found}.
    \item \textbf{No matching files}: When criteria are provided but no file satisfies them,
          \texttt{No file found} is printed.
    \item \textbf{Permission denied}: Directories that cannot be opened (e.g. \texttt{/etc/ssl/private})
          print an error to stderr and traversal continues with the next entry.
    \item \textbf{Deep nesting}: The program correctly indents entries at all depth levels
          using the \texttt{HYPHEN\_PER\_LEVEL} constant.
    \item \textbf{Multiple \texttt{+} operators}: Patterns like \texttt{pa+s+wd} correctly
          match strings with repeated characters at multiple positions.
\end{itemize}

% ─── 7. Challenges ────────────────────────────────────────────────────────────
\section{Challenges and Solutions}

\subsection{Implementing the \texttt{+} Regex Operator}

The main challenge was correctly handling the \texttt{+} operator with an independent
index variable (\texttt{k}) that tracks the current position in the string separately
from the pattern index (\texttt{j}). An early version lacked bounds checking in the
\texttt{while} loop consuming repeated characters, which caused out-of-bounds reads.
The fix was adding the guard \texttt{i + k < len\_of\_string} before accessing
\texttt{string[i + k]}.

\subsection{Missing \texttt{if} Guard for File Type Check}

During development, the \texttt{if(criteria->file\_type != '\textbackslash0')} guard was
accidentally omitted in \texttt{check\_matches()}. This caused \texttt{check\_type()} to
always run with \texttt{'\textbackslash0'} as the expected type, making every file fail the
check — even when \texttt{-t} was not provided. This was discovered by isolating
\texttt{check\_matches()} in a separate test program.

\subsection{Choosing Between \texttt{exit()} and \texttt{volatile} Flag for SIGINT}

Two approaches were considered for SIGINT handling: calling \texttt{exit()} directly from
the signal handler, or setting a \texttt{volatile sig\_atomic\_t} flag and checking it in
the traversal loop. The flag approach was chosen because it allows the program to close
open directory handles cleanly before exiting, following POSIX best practices for
async-signal-safe code.

\subsection{Using \texttt{lstat()} vs \texttt{stat()}}

An early version used \texttt{stat()}, which caused symbolic links to be misidentified
as their target file type. Switching to \texttt{lstat()} resolved this and correctly
identifies symbolic links when using the \texttt{-t l} criterion.

% ─── 8. Conclusion ────────────────────────────────────────────────────────────
\section{Conclusion}

This homework required the design and implementation of a fully functional, POSIX-compatible
recursive file search program in C. The program was developed with a strong emphasis on
modularity, correctness, and robustness under real system conditions.

Each module was designed with a single, well-defined responsibility, making the codebase
easy to follow and maintain. The custom \texttt{+} regex operator was implemented entirely
from scratch without any external library, deepening the understanding of string matching
algorithms. The SIGINT handler was implemented using \texttt{sigaction()} and a
\texttt{volatile sig\_atomic\_t} flag rather than a simple \texttt{exit()} call, reflecting
POSIX best practices for signal-safe code.

Throughout the development process, several real bugs were encountered and resolved —
including a missing sentinel guard for the file type criterion and a buffer overflow risk
in the regex engine — which reinforced the importance of careful testing and isolation
of individual components.

The final program passes all test cases, produces zero Valgrind errors, handles CTRL-C
gracefully, and compiles cleanly with \texttt{-Wall -Wextra} without any warnings.
All requirements specified in the homework description have been satisfied.

% ─── 9. References ────────────────────────────────────────────────────────────
\section{References}

\begin{enumerate}[noitemsep]

    \item Michael Kerrisk. \textit{The Linux Programming Interface}. No Starch Press, 2010.

    \item Gebze Technical University. \textit{CSE344 System Programming -- Homework \#1 Assignment Sheet}. Spring 2026.

    \item Erkan Zergeroğlu. \textit{CSE344 System Programming -- Course Lecture Notes, Slides and Course Materials}. Gebze Technical University, Spring 2026.

    \item Anthropic. \textit{Claude AI Assistant} (claude.ai) --- Used as a learning and reference aid during development.

    \item Google. \textit{Gemini AI Assistant} (gemini.google.com) --- Used as a learning and reference aid during development.

\end{enumerate}

\end{document}